\documentclass[12pt, a4paper]{ctexart}
\usepackage{geometry}
\geometry{left=2.5cm, right=2.5cm, top=3cm, bottom=3cm}
\usepackage{amsmath}
\usepackage{circuitikz}
\usepackage{float}
\usepackage{multirow}
\usepackage{arydshln}
\usetikzlibrary{arrows, shapes.gates.logic.US}
\ctikzset{
    logic ports=ieee,
    logic ports/scale=0.8,
    multipoles/flipflop/pin spacing=0.4
}

\begin{document}

\section*{第6章习题}

\subsection*{3.}

\begin{itemize}
    \item \texttt{func1} 功能: 保留原始数值的最低 8 位, 其他位补零。
    \item \texttt{func2} 功能: 保留原始数值的最低 8 位, 其他位补符号位。
\end{itemize}

\begin{table}[H]
    \centering
    \begin{tabular}{|c|c|c|c|c|c|}
        \hline
        \multicolumn{2}{|c|}{\texttt{w}} & \multicolumn{2}{c|}{\texttt{func1(w)}} & \multicolumn{2}{c|}{\texttt{func2(w)}} \\
        \hline
        机器数 & 值 & 机器数 & 值 & 机器数 & 值 \\
        \hline
        \texttt{0x0000007F} & 127 & \texttt{0x0000007F} & 127 & \texttt{0x0000007F} &  127 \\
        \texttt{0x00000080} & 128 & \texttt{0x00000080} & 128 & \texttt{0xFFFFFF80} & -128 \\
        \texttt{0x000000FF} & 255 & \texttt{0x000000FF} & 255 & \texttt{0xFFFFFFFF} &   -1 \\
        \texttt{0x00000100} & 256 & \texttt{0x00000000} &   0 & \texttt{0x00000000} &    0 \\
        \hline
    \end{tabular}
    \caption{习题 3 用表}
\end{table}

\subsection*{4.}

\begin{table}[H]
    \centering
    \begin{tabular}{|c|c|c|c|c|c|c|c|c|}
        \hline
        \multirow{2}{*}{模式} & \multicolumn{2}{c|}{$x$} & \multicolumn{2}{c|}{$y$} & \multicolumn{2}{c|}{$x\times y$(截断前)} & \multicolumn{2}{c|}{$x\times y$(截断后)} \\
        \cline{2-9}
        & 机器数 & 值 & 机器数 & 值 & 机器数 & 值 & 机器数 & 值 \\
        \hline
        无符号数 & 110 & 6 & 010 & 2 & 001100 & 12 & 100 & 4 \\
        二进制补码 & 110 & -2 & 010 & 2 & 111100 & -4 & 100 & -4 \\
        \hline
        无符号数 & 001 & 1 & 111 & 7 & 000111 & 7 & 111 & 7 \\
        二进制补码 & 001 & 1 & 111 & -1 & 111111 & -1 & 111 & -1 \\
        \hline
        无符号数 & 111 & 7 & 111 & 7 & 110001 & 49 & 001 & 1 \\
        二进制补码 & 111 & -1 & 111 & -1 & 000001 & 1 & 001 & 1 \\
        \hline
    \end{tabular}
    \caption{习题 4 用表}
\end{table}

\subsection*{5.}

$M=15,N=4$

\subsection*{6.}

记 $G_i=A_iB_i,P_i=A_i+B_i$，则有

\subsubsection*{(1) 串行进位}

\begin{itemize}
    \item $C_1=G_1+P_1C_0$
    \item $C_2=G_2+P_2C_1$
    \item $C_3=G_3+P_3C_2$
    \item $C_4=G_4+P_4C_3$
\end{itemize}

\subsubsection*{(2) 先行进位}

\begin{itemize}
    \item $C_1=G_1+P_1C_0$
    \item $C_2=G_2+P_2G_1+P_2P_1C_0$
    \item $C_3=G_3+P_3G_2+P_3P_2G_1+P_3P_2P_1C_0$
    \item $C_4=G_4+P_4G_3+P_4P_3G_2+P_4P_3P_2G_1+P_4P_3P_2P_1C_0$
\end{itemize}

\subsection*{7.}

\subsubsection*{(1)}

\begin{equation}
    \left[x+y\right]_\text{补}=0010B=2,
    \left[x-y\right]_\text{补}=1000B=-8
\end{equation}

\subsubsection*{(2)}

\begin{table}[H]
    \centering
    \begin{tabular}{cccccccc}
        \hline
        \multicolumn{4}{c}{P} && \multicolumn{3}{c}{Y} \\
        \hline
          & 0&0&0 &\multicolumn{1}{c|}{}& 1&0&1 \\ \hdashline
        + & 1&0&1 &\multicolumn{1}{c|}{}&       \\ \cline{1-4}
        0 & 1&0&1 &\multicolumn{1}{c|}{}& 1&0&1 \\ \cline{6-6}
          & 0&1&0 && \multicolumn{1}{c|}{1}&1&0 \\ \hdashline
        + & 0&0&0 && \multicolumn{1}{c|}{}      \\ \cline{1-4}
        0 & 0&1&0 && \multicolumn{1}{c|}{1}&1&0 \\ \cline{7-7}
          & 0&0&1 && 0&\multicolumn{1}{c|}{1}&1 \\ \hdashline
        + & 1&0&1 && \multicolumn{2}{c|}{}      \\ \cline{1-4}
        0 & 1&1&0 && 0&\multicolumn{1}{c|}{1}&1 \\ \cline{8-8}
          & 0&1&1 && 0&0&\multicolumn{1}{c|}{1} \\ \hline
    \end{tabular}
\end{table}

所以 $\left[x\times y\right]_\text{原}=1011001B=-25$

\subsubsection*{(3)}

\begin{table}[H]
    \centering
    \begin{tabular}{ccccccccccc}
        \hline
        & \multicolumn{4}{c}{P} && \multicolumn{4}{c}{Y} & Y$_{-1}$ \\
        \hline
          & 0&0&0&0 &\multicolumn{1}{c|}{}& 1&1&0&1 & 0 \\ \hdashline
        + & 0&1&0&1 &\multicolumn{1}{c|}{}              \\ \cline{1-5}
        0 & 0&1&0&1 &\multicolumn{1}{c|}{}              \\ \cline{7-8}
          & 0&0&0&1 && 0&\multicolumn{1}{c|}{1}&1&1 & 0 \\ \hdashline
        + & 1&0&1&1 && &\multicolumn{1}{c|}{}           \\ \cline{1-5}
        1 & 1&1&0&0 && &\multicolumn{1}{c|}{}           \\ \cline{9-10}
          & 1&1&1&1 && 0&0&0&\multicolumn{1}{c|}{1}     \\ \hline
    \end{tabular}
\end{table}

所以 $\left[x\times y\right]_\text{补}=11110011B=-15$

\subsubsection*{(4)}

$X=0101,Y=0101,-Y=1011$，有

\begin{table}[H]
    \centering
    \begin{tabular}{cccccccccc}
        \hline
        & \multicolumn{4}{c}{R} && \multicolumn{4}{c}{Q} \\
        \hline
          & 0&0&0&0 && 1&0&\multicolumn{1}{c|}{1}   \\ \hdashline
        + & 1&0&1&1 && \multicolumn{3}{c|}{}        \\ \cline{1-5}
          & 1&0&1&1 && 1&0&\multicolumn{1}{c|}{1}&0 \\ \hdashline
          & 0&1&1&1 && 0&\multicolumn{1}{c|}{1}&0   \\
        + & 0&1&0&1 && \multicolumn{2}{c|}{}        \\ \cline{1-5}
          & 1&1&0&0 && 0&\multicolumn{1}{c|}{1}&0&0 \\ \hdashline
          & 1&0&0&0 && \multicolumn{1}{c|}{1}&0&0   \\
        + & 0&1&0&1 && \multicolumn{1}{c|}{}        \\ \cline{1-5}
          & 1&1&0&1 && \multicolumn{1}{c|}{1}&0&0&0 \\ \hdashline
          & 1&0&1&1 &\multicolumn{1}{c|}{}& 0&0&0   \\
        + & 0&1&0&1 &\multicolumn{1}{c|}{}          \\ \cline{1-5}
          & 0&0&0&0 &\multicolumn{1}{c|}{}& 0&0&0&1 \\ \hline
    \end{tabular}
\end{table}

所以 $\left[x/y\right]_\text{原}=-0001B=-1$，余数为 $0000B=0$

\subsubsection*{(5)}

$X=00000101,Y=1101,-Y=0011$，有

\begin{table}[H]
    \centering
    \begin{tabular}{cccccccccc}
        \hline
        & \multicolumn{4}{c}{R} && \multicolumn{4}{c}{Q} \\
        \hline
          & 0&0&0&0 && 0&1&0&\multicolumn{1}{c|}{1} \\ \hdashline
        + & 1&1&0&1 &\multicolumn{5}{c|}{}          \\ \cline{1-5}
          & 1&1&0&1 && 0&1&0&\multicolumn{1}{c|}{1} \\ \hdashline
          & 1&0&1&0 && 1&0&\multicolumn{1}{c|}{1}&1 \\
        + & 0&0&1&1 &\multicolumn{4}{c|}{}          \\ \cline{1-5}
          & 1&1&0&1 && 1&0&\multicolumn{1}{c|}{1}&1 \\ \hdashline
          & 1&0&1&1 && 0&\multicolumn{1}{c|}{1}&1&1 \\
        + & 0&0&1&1 &\multicolumn{3}{c|}{}          \\ \cline{1-5}
          & 1&1&1&0 && 0&\multicolumn{1}{c|}{1}&1&1 \\ \hdashline
          & 1&1&0&0 && \multicolumn{1}{c|}{1}&1&1&1 \\
        + & 0&0&1&1 &\multicolumn{2}{c|}{}          \\ \cline{1-5}
          & 1&1&1&1 && \multicolumn{1}{c|}{1}&1&1&1 \\ \hdashline
          & 1&1&1&1 &\multicolumn{1}{c|}{}& 1&1&1&1 \\
        + & 0&0&1&1 &\multicolumn{1}{c|}{}          \\ \cline{1-5}
          & 0&0&1&0 &\multicolumn{1}{c|}{}& 1&1&1&0 \\
          & &&&     && + &&&1                       \\ \cline{7-10}
          & 0&0&1&0 && 1&1&1&1                      \\ \hline
    \end{tabular}
\end{table}

所以 $\left[x/y\right]_\text{补}=1111B=-1$，余数为 $0010B=2$

\subsection*{9.}

\begin{itemize}
  \item 左规: 运算结果为 $0.b\cdots b$ 时，需要进行左规。
  将尾数的数值位逐次左移，尾数每左移一位，对应的阶码减 1，直到将第一位“1”移到小数点左边，使其变成标准的规格化形式。
  \item 右规: 运算结果为 $1b.b\cdots b$ 时，需要进行右规。
  将尾数整体右移一位，同时对应的阶码加 1，使尾数恢复为规格化形式。
\end{itemize}

\subsection*{12.}

\subsubsection*{(1)}

由
$x=0.75=1.1B\times2^{-1},y=-65.25=-1.00000101B\times2^6$
有
$s_x=0,E_x=01111110B,M_x=0(1).10\cdots0,
s_y=1,E_y=10000101B,M_y=1(1).000001010\cdots0$
，所以
$\Delta E=-7,M'_x=0(0).000000110\cdots0$
，所以
$M_b=1(1).00000010\cdots0,E_b=10000101B$
，所以
$0.75+(-65.25)=-1.0000001B\times2^6=-64.5$

\end{document}